\documentclass{article}
\usepackage{amsmath,amssymb}
\usepackage{xurl}
\usepackage[hidelinks]{hyperref}
% Shared commands for notes in docs/paper-gaps/.

\DeclareUnicodeCharacter{2080}{\ifmmode{}_{0}\else\textsubscript{0}\fi}
\DeclareUnicodeCharacter{2081}{\ifmmode{}_{1}\else\textsubscript{1}\fi}
\DeclareUnicodeCharacter{2082}{\ifmmode{}_{2}\else\textsubscript{2}\fi}
\DeclareUnicodeCharacter{2099}{\ifmmode{}_{n}\else\textsubscript{n}\fi}

\newcommand{\Lean}{\textsc{Lean}}
\ExplSyntaxOn
\NewDocumentCommand{\leanid}{m}{
  \ifmmode
    \textsf{\detokenize{#1}}
  \else
    \group_begin:
    \urlstyle{sf}
    \tl_set:Nn \l_tmpa_tl {#1}
    \tl_replace_all:Nnn \l_tmpa_tl {\_} {_}
    \exp_args:NV \path \l_tmpa_tl
    \group_end:
  \fi
}
\ExplSyntaxOff
\providecommand{\tr}{\operatorname{tr}}
\newcommand{\tnleanIssueBase}{https://github.com/LionSR/TNLean/issues/}
\newcommand{\tnleanPrBase}{https://github.com/LionSR/TNLean/pull/}
\newcommand{\tnleanDocBase}{https://sirui-lu.com/TNLean/docs/}
\newcommand{\ghissue}[1]{\href{\tnleanIssueBase#1}{GitHub issue \##1}}
\newcommand{\ghpr}[1]{\href{\tnleanPrBase#1}{PR \##1}}
\newcommand{\doclink}[1]{\href{\tnleanDocBase#1.html}{\path{#1}}}

% Dirac notation and the identity operator, as in blueprint/src/macros/common.tex.
% \providecommand keeps note-local definitions authoritative.
\usepackage{dsfont}
\providecommand{\ket}[1]{|#1\rangle}
\providecommand{\bra}[1]{\langle #1|}
\providecommand{\braket}[2]{\langle #1 | #2 \rangle}
\providecommand{\ketbra}[2]{|#1\rangle\!\langle #2|}
\providecommand{\Id}{\mathds{1}}
\providecommand{\one}{\mathds{1}}
\providecommand{\C}{\mathbb{C}}
\providecommand{\MN}[1]{M_{#1}(\C)}
\providecommand{\MD}{M_D(\C)}

% Verdict marker: \gapnote{<kind>}{<status>}. Kind is the policy
% classification (clarification, local-correction, scope-restriction,
% unfaithful, false-source, open-gap); status is open, wip (elimination
% underway), resolved, or historical. Machine-read by the site index;
% prints a verdict line.
\newcommand{\gapnote}[2]{\par\noindent\textbf{Verdict:} #1 (#2).\par}

\title{Normalization and Gate Orientation in the MPU--QCA Index Comparison}
\author{}
\date{2026-10-29}
\begin{document}
\maketitle
\gapnote{local-correction}{open}

\section{Two printed normalizations}

Definition~\texttt{def:index} of Ref.~\cite{Cirac2017MPU}, source
\path{Papers/1703.09188/paper_v2.tex}, lines 681--688, defines the MPU index by
\begin{align}
    \operatorname{ind}_{\mathrm{MPU}}
    =\tfrac12(\log_2r-\log_2\ell)
    =\log_2(r/d)=-\log_2(\ell/d),
    \qquad r\ell=d^2.
    \label{eq:mpu_qca_main_index_normalization}
\end{align}
The last equality is also stated explicitly at line 688.

In the Appendix, equation~\texttt{eq:original-index}, lines 2329--2333,
identifies the positive rational QCA index as
\begin{align}
    \operatorname{ind}_{\mathrm{QCA}}=s/d=\ell/d.
    \label{eq:mpu_qca_original_ratio}
\end{align}
Line 2334 instead prints
$\operatorname{ind}_{\mathrm{MPU}}=-2\log_2(\ell/d)$.
This factor of two conflicts with both the defining formula and its
immediate consequence in (\ref{eq:mpu_qca_main_index_normalization}).

\section{Nondegenerate check and correction}

For $d>1$, $r=d^2$ and $\ell=1$, the defining formula gives
$\log_2d$, whereas the printed Appendix expression gives $2\log_2d$.
These are the nonzero shift dimensions, not an empty or inactive factor.
The correction is therefore
\begin{align}
    \operatorname{ind}_{\mathrm{MPU}}
    =-\log_2(\operatorname{ind}_{\mathrm{QCA}}).
    \label{eq:mpu_qca_corrected_log_relation}
\end{align}
The positive rational identification (\ref{eq:mpu_qca_original_ratio}) is
unchanged. The left support factor carries $\ell$, as printed at line 2333.
Main-text line 543 states that $u$ and $v^\dagger$ map the physical pair
space to the intermediate $r\ell$-dimensional space. Thus $v$ maps back
to the physical pair, in the direction used by the actual forward
conjugation in equation~\texttt{eq:appendix-1}. Transporting the intermediate
left factor to the physical pair therefore gives
$v(I_r\otimes M_\ell)v^\dagger$. Appendix line 2333 instead writes
$v^\dagger(I\otimes M_\ell)v$, reversing the adjoints relative to the
main-text direction. The formalized finite output family uses the actual
forward order; its least left support is derived, rather than assumed
from either printed expression. The matrix size is $\ell$ in the corrected
presentation, so this adjoint-order correction does not change the
positive rational index.

\section{Formalization scope}

The finite output-family calculation derives its least left support
algebra and matrix size from the supplied unitary gates. It does not
assume a support-factor presentation or identify an arbitrary supplied
gate with a selected singular-value factorization. The remaining actual
QCA comparison must identify this finite family with the canonical
nearest-neighbour pair image of the induced automorphism before applying
the GNVW index construction. That assembly is distinct from the printed
factor-of-two correction, and remains pending Lean checking.

\bibliographystyle{alpha}
\bibliography{references}
\end{document}
